\documentclass{article}
\usepackage{amsmath, amssymb, amsthm}
\usepackage[margin=1in]{geometry}

\begin{document}

\begin{center}
    \Large\textbf{MHF 3302 (Test 2, Fall 2023) Section 1} \\
    \vspace{0.2cm}
    \normalsize Name \underline{\hspace{3cm} KEY \hspace{3cm}}
\end{center}

\vspace{0.5cm}
\noindent Show your work in the space allocated.
\vspace{0.5cm}

\section*{Question 1}
\textbf{15 pts. \textcircled{1}} Define $P(n): 2 + 2^2 + 2^3 + \dots + 2^{n-1} = 2^n - 2, \; n \ge 2$. \\
Use PMI to prove $P(n)$ is true $\forall n \ge 2$.

\begin{itemize}
    \item[(i)] $P(2): 2 \overset{?}{=} 2^2 - 2$ (true)
    \item[(ii)] Assume $P(k)$ true, ie $2 + 2^2 + 2^3 + \dots + 2^{k-1} = 2^k - 2$
    
    $(2 + 2^2 + 2^3 + \dots + 2^{k-1}) + 2^k = 2^k - 2 + 2^k = 2 \cdot 2^k - 2 = 2^{k+1} - 2$
    
    $\therefore P(k+1)$ is true.
    
    Hence $P(n)$ is true $\forall n \ge 2$.
\end{itemize}

\vspace{0.3cm}
\noindent \textcircled{1}' Prove $P(n): 1+2+\dots+n = \frac{n(n+1)}{2}, n \ge 1$ \\
\noindent \textcircled{1}'' Prove $P(n): 8 | (9^n-1), n \ge 0$

\vspace{0.5cm}

\section*{Question 2}
\textbf{20 pts. \textcircled{2}} (i) Fill in the table below for multiplication in $\mathbb{Z}_3$:

\begin{center}
$
\begin{array}{c|ccc}
\cdot & [0] & [1] & [2] \\
\hline
[0] & [0] & [0] & [0] \\
[1] & [0] & [1] & [2] \\
[2] & [0] & [2] & [1]
\end{array}
$
\end{center}

\vspace{0.3cm}
\noindent The multiplicative inverse of $[2]$ in $\mathbb{Z}_3$ is \underline{$[2]$}. Also, $47 \equiv b \pmod 3$ implies that $b = \underline{\; 2 \;}$, where $0 \le b < 3$. \\
$3 | (47-2) \text{ or } 3 | 45$

\vspace{0.5cm}

\noindent (ii) Let $x, y \in \mathbb{Z}$. Use the Direct Method to prove: \\
$x \equiv 2 \pmod 4 \text{ and } y \equiv 1 \pmod 4 \implies 4 \nmid (x^2 + 2y)$.

\vspace{0.3cm}
\noindent $x \equiv 2 \pmod 4 \text{ and } y \equiv 1 \pmod 4 \implies 4 \nmid (x^2+2y)$. \\
$x - 2 = 4q, \; y - 1 = 4p \quad \exists q, p \in \mathbb{Z}$. Then \\
$x^2 + 2y = (4q+2)^2 + 2(4p+1) = 16q^2 + 16q + 4 + 8p + 2$ \\
$= 4\underbrace{(4q^2 + 4q + 2p + 1)}_{\in \mathbb{Z}} + 2 \quad \therefore 4 \nmid (x^2 + 2y)$

\vspace{0.5cm}

\section*{Question 3}
\textbf{15 pts. \textcircled{3}} 10 pts (i) Prove that $\sqrt{2} \in \mathbb{I}$ (irrational).

\vspace{0.3cm}
\noindent Pf. Assume that $\sqrt{2} = p/q, \; p, q \in \mathbb{Z}, \; q \neq 0$, and $p/q$ reduced to lowest terms. Then $2q^2 = p^2 \implies 2 | p^2 \implies 2 | p \implies p = 2r \;\exists r \in \mathbb{Z}$. Then $2q^2 = 4r^2 \text{ or } q^2 = 2r^2 \implies 2 | q^2 \implies 2 | q$. Hence $2 | p \text{ and } 2 | q$, a contradiction. \\
$\therefore \sqrt{2} \in \mathbb{I}$.

\vspace{0.5cm}
\noindent 5 pts. (ii) Give an irrational number between 10 and 11: \\
\underline{\quad $10 + \frac{\sqrt{2}}{2}$ \quad}.

\vspace{0.5cm}

\section*{Question 4}
\textbf{15 pts. \textcircled{4}} Define statement $S: \forall a \in \mathbb{Z}, \forall b \in \mathbb{Z}, a \text{ odd and } b \text{ odd } \implies 4 \nmid (3a^2 + 7b^2)$

\vspace{0.3cm}
\noindent 5 pts. (i) $\sim S$: \underline{ $\exists a \in \mathbb{Z}, \exists b \in \mathbb{Z}, a \text{ odd and } b \text{ odd}$ } \\
\hspace{2.5cm} \underline{ $\text{and } 4 | (3a^2 + 7b^2)$ }

\vspace{0.3cm}
\noindent 10 pts. (ii) Assume that $\sim S$ is true and prove that this leads to a contradiction.

\vspace{0.3cm}
\noindent Assume that $\sim S$ is true. Then $a = 2x+1 \text{ and } b = 2y+1 \;\exists x, y \in \mathbb{Z}$. Hence $3a^2 + 7b^2 = 3(2x+1)^2 + 7(2y+1)^2$ \\
$= 12x^2 + 12x + 3 + 28y^2 + 28y + 7$ \\
$= 4\underbrace{[3x^2 + 3x + 7y^2 + 7y + 2]}_{\in \mathbb{Z}} + 2$, a contradiction. \\
$\therefore S \text{ is true}$

\vspace{0.5cm}

\section*{Question 5}
\textbf{15 pts. \textcircled{5}} (i) Complete the Intermediate Value THM: \\
Let $f: [a,b] \to \mathbb{R}$. If $f$ satisfies: \\
(a) \underline{ $f \text{ is continuous on } [a,b]$ } \\
(b) \underline{ $f(a) < 0 \text{ and } f(b) > 0 \text{ or vice-versa}$ } \\
Then \underline{ $\exists x_0 \in [a,b] \text{ such that } f(x_0) = 0$. }

\vspace{0.3cm}
\noindent (ii) Use (i) to prove that $x^6 + x^4 + 2x^2 - 2 = 0$ has at least one solution $x_0$ in the interval $[0,1]$.

\vspace{0.3cm}
\noindent Choose $f(x) = x^6 + x^4 + 2x^2 - 2$. Then $f$ is a polynomial and thus continuous. Further, $f(0) = -2 < 0 \text{ and } f(1) = 2 > 0$. Therefore by IVT $\exists x_0 \in [0,1]$ such that $f(x_0) = 0$, \\
i.e., $x_0^6 + x_0^4 + 2x_0^2 - 2 = 0$.

\vspace{0.3cm}
\noindent (iii) Assume that $x_0$ and $x_1$ are solutions of (ii) above. If $x_0 < x_1$, arrive at a contradiction.

\vspace{0.3cm}
\noindent Suppose that $x_0 < x_1 \text{ and } x_0^6 + x_0^4 + 2x_0^2 - 2 = 0 = x_1^6 + x_1^4 + 2x_1^2 - 2$. Then $(x_0^6 - x_1^6) + (x_0^4 - x_1^4) + 2(x_0^2 - x_1^2) = 0$, a contradiction since sum is $< 0$. \\
$\therefore x_0 = x_1 \text{ and solution is unique.}$

\vspace{0.5cm}

\section*{Question 6}
\textbf{20 pts. \textcircled{6}} Define $S: \forall x \in \mathbb{R}, \forall y \in \mathbb{R}, 3 | x \text{ and } 3 | (x+y) \implies 3 | y$. \textit{(Note: Domain intended is $\mathbb{Z}$ based on divisibility operators.)}

\vspace{0.3cm}
\noindent (i) Write the Contrapositive of $S$: \underline{ $\forall x \in \mathbb{R}, \forall y \in \mathbb{R},$ } \\
\hspace{0.5cm} \underline{ $3 \nmid y \implies 3 \nmid x \text{ or } 3 \nmid (x+y)$ }.

\vspace{0.3cm}
\noindent (ii) Write $\sim S$: \underline{ $\exists x \in \mathbb{R}, \exists y \in \mathbb{R}, 3 | x \text{ and } 3 | (x+y)$ } \\
\hspace{0.5cm} \underline{ $\text{and } 3 \nmid y$. }

\vspace{0.3cm}
\noindent (iii) Using the Direct Method, prove $S$ is true.

\vspace{0.3cm}
\noindent Pf. $x = 3p, \; x+y = 3q \quad \exists p, q \in \mathbb{Z}$. Then $y = 3q - x = 3q - 3p = 3(q-p) \quad \therefore 3 | y$.

\vspace{0.3cm}
\noindent (iv) Prove $S$ is true using a Contradiction Argument.

\vspace{0.3cm}
\noindent Pf. Assume that $\sim S$ is true. Then $x = 3p \text{ and } x+y = 3q \quad \exists p, q \in \mathbb{Z}$. Then $y = 3q - x = 3q - 3p = 3(q-p) \text{ and thus } 3 | y$, a contradiction. $\therefore S \text{ is true}$.

\end{document}